\documentclass[11pt]{article}
\usepackage[T1]{fontenc}
\usepackage[utf8]{inputenc}
\usepackage{lmodern}
\usepackage{amsmath,amssymb,amsthm,mathtools}
\usepackage{booktabs,array,longtable,tabularx}
\usepackage[margin=0.78in]{geometry}
\usepackage[hidelinks]{hyperref}
\usepackage{xcolor}
\usepackage{enumitem}
\usepackage{microtype}
\usepackage{url}

\newtheorem{theorem}{Theorem}
\newtheorem{proposition}{Proposition}
\newtheorem{lemma}{Lemma}
\newtheorem{corollary}{Corollary}
\newcommand{\K}{\mathbb{K}}
\newcommand{\Q}{\mathbb{Q}}
\newcommand{\R}{\mathbb{R}}
\newcommand{\Xsrc}{X_{\mathrm{src}}}
\newcommand{\Rsrc}{R_{\mathrm{src}}}
\newcommand{\Usrc}{U_{\mathrm{src}}}
\newcommand{\bsrc}{b_{\mathrm{src}}}
\newcommand{\Osrc}{O_{\mathrm{src}}}
\newcommand{\Bridge}{\operatorname{Bridge}_{\mathrm{OM}}}
\newcommand{\Viab}{\operatorname{Viab}}
\newcommand{\status}[1]{\textsf{#1}}
\newcolumntype{P}[1]{>{\raggedright\arraybackslash}p{#1}}
\setlist{nosep,leftmargin=*}
\urlstyle{tt}

\title{Post-Oriented-Matroid Refuter Route Selection\\
\large V22 P4--T02 B2 theoretical-continuation packet}
\author{Theoretical Continuation Selector parent synthesis, RT-20260811-002}
\date{11 August 2026}

\begin{document}
\maketitle

\begin{center}
\fbox{\begin{minipage}{0.94\linewidth}
\textbf{Decisive result.} Exactly one future packet is selected but not
executed:
\begin{center}
\small\path{PKT-V22-P4T02-B2-SOURCE-DYNAMICAL-VIABILITY-ADMISSIBILITY-SELECTOR-FORMALIZATION-V1}
\end{center}
under \path{ontology-formalizer@0.2.0}.  It is an
\status{ontology-law-research-packet} triggered by
\status{derivation-critical missing source law}.  The packet must formalize,
from target-free source premises, an occurrence or component selector together
with an admissible-transition and viability law, or return the smallest exact
missing primitive or countermodel.

\smallskip
\textbf{Boundary.} The selection is \status{draft/control},
\status{proposal-only}, and \status{source-extension data}.  It does not
construct, prove, audit, stress, or adopt the selected law.  It preserves the
RT018 scoped obstruction, all four local freezes, D7 non-evaluation, B2
inactivity, the P4--T03 lock, and all downstream physics and protected gates.
Every Distance-to-GR burden remains unchanged.
\end{minipage}}
\end{center}

\section{Fixed predecessor and selection question}

RT018 preserves the arbitrary-finite-rank oriented-matroid classifier on its
declared exact domain while establishing a candidate-local obstruction:
pointwise totality, presentation equivariance, refinement bookkeeping, and
wall enumeration do not themselves select which source realization occurs,
which variations are admissible, or which branch is physically meaningful.
The exact obstruction is
\begin{center}
\small\path{OB-V22-P4T02-B2-OM-BRIDGE-SELECTION-ROBUSTNESS-001}.
\end{center}
Diagnostic use of $\Bridge$ remains open.  The unchanged standalone bridge is
locally frozen by
\begin{center}
\small\path{NDCL-V22-P4T02-B2-ORIENTED-MATROID-BRIDGE-SELECTION-ROBUSTNESS}.
\end{center}
Three earlier local freezes also remain active:
\begin{itemize}
  \item \path{NDCL-V22-P4T02-B2-SHARED-TAU-SELECTOR};
  \item \path{NDCL-V22-P4T02-B2-REPAIRED-QUOTIENT-DESCENT-ROBUSTNESS};
  \item \path{NDCL-V22-P4T02-B2-COMMON-CHARACTER-HOLONOMY-PROTECTION}.
\end{itemize}

The present task compares exactly four continuations named by handoff-1015:
a source-side selection or admissibility theorem, a realization-sensitive
bridge, a source-side irrelevance theorem, and a protected human-gated
ontology stop.  It may choose one future packet; it may execute none.

\section{Selection and protection are distinct structures}

Let $\Xsrc$ be a declared full source-state or source-realization space.  The
future packet may use a source successor relation
\[
  \Rsrc\subseteq \Xsrc\times\Xsrc
\]
or a source semiflow, but its provenance must be derived from tracked source
equations or explicitly classified as proposal-only source-extension data.
Let $\Usrc\subseteq\Xsrc$ be a source-admissible domain fixed before desired
physical behavior is inspected.  An occurrence law $\Osrc$ selects an initial
state or a unique source component.  A viability law restricts admitted
successors so trajectories remain in $\Usrc$.

\begin{theorem}[Occurrence--viability independence]
A viability law without an occurrence or component selector protects supplied
initial states but selects none.  An occurrence selector without a viability
law chooses initial data but does not protect the selected state from an
admitted trajectory leaving $\Usrc$.
\end{theorem}

\begin{proof}
The viability statement has the form ``for every supplied $x_0$ satisfying a
predicate, every admitted successor trajectory remains in the domain.''  It
contains no existence or uniqueness clause selecting one $x_0$ or component.
Conversely, an occurrence statement constrains $x_0$ but contains no universal
condition on successors.  A model with two invariant admissible components
and no occurrence law satisfies viability but not selection.  A model with a
unique selected $x_0$ and one admitted successor outside $\Usrc$ satisfies
occurrence but not viability.  Thus neither structure entails the other.
\end{proof}

This logical separation is the key route discriminator.  A quantitative
realization margin can measure a supplied state's distance from a wall, but it
does not by itself choose which realization occurs.  A pointwise classifier
can label every supplied state, but it supplies neither clause.

\begin{lemma}[Representation--occurrence separation]
Even an injective representation $J:X\to Y$ of a source-state set with at
least two elements supplies neither a distinguished element of $X$, an
occurrence map into $X$, nor an admissible transition relation on $X$.
\end{lemma}

\begin{proof}
Take $X=\{x_0,x_1\}$ and $J=\operatorname{id}_X$.  The same injective map is
compatible with occurrence of $x_0$, occurrence of $x_1$, or no selected
occurrence, and it contains no subset of $X\times X$ declaring allowed
changes.  Injectivity distinguishes supplied states; it does not cause,
select, or evolve them.
\end{proof}

Let $\beta:\Xsrc\to B$ be any total diagnostic classifier and let
$\operatorname{Reach}_{\Rsrc}(x)$ contain the endpoints of all finite
$\Rsrc$-paths beginning at $x$, including the length-zero path.

\begin{theorem}[Reachability-component branch robustness]
The diagnostic branch is preserved along every finite admitted path beginning
at $x$ if and only if $\beta$ is constant on
$\operatorname{Reach}_{\Rsrc}(x)$.  Equivalently, there is an
$\Rsrc$-forward-invariant set $U$ containing $x$ on which $\beta$ is
constant.  When robustness holds,
$\operatorname{Reach}_{\Rsrc}(x)$ is the least such forward-invariant carrier.
\end{theorem}

\begin{proof}
The reachable set contains $x$ and is forward invariant by path
concatenation.  If every admitted path preserves the branch, each reachable
endpoint $y$ satisfies $\beta(y)=\beta(x)$, so $\beta$ is constant on the
reachable set.  The converse is immediate because every finite path endpoint
is reachable.  If a forward-invariant $U$ contains $x$, induction on path
length gives $\operatorname{Reach}_{\Rsrc}(x)\subseteq U$; taking the reachable
set itself proves leastness and the equivalent formulation.
\end{proof}

For RT018's $A(t)=[e_1,e_2,(-1,t)]$, any admitted relation connecting
$A(1/n)$ to $A(-1/n)$ has a branch-inhomogeneous reachable set.  Pointwise
totality cannot repair that failure.  Conversely, forbidding the connection
only after inspecting $\beta$ is goal-property preload.  A valid repair must
derive both the occurrence state and the transition relation independently
from source premises.

\section{A conditional source-barrier theorem}

The selected future packet is required to seek an exact source-side barrier
or to explain precisely why no such representation exists on its declared
domain.

\begin{theorem}[Finite source-barrier propagation]
Let $\K$ be an ordered field, let $\bsrc:\Xsrc\to\K$, and define
$\Usrc=\{x:\bsrc(x)>0\}$.  Suppose one $\lambda>0$ satisfies
\[
  \bsrc(y)\geq \lambda\,\bsrc(x)
  \qquad\text{for every admitted successor }y\in\Rsrc(x).
\]
Then every finite admitted trajectory $x_0,\ldots,x_n$ beginning in $\Usrc$
remains in $\Usrc$, with
\[
  \bsrc(x_n)\geq \lambda^n\bsrc(x_0)>0.
\]
\end{theorem}

\begin{proof}
Induct on $n$.  The case $n=0$ is the hypothesis.  If the bound holds at
$n$, the successor inequality gives
$\bsrc(x_{n+1})\geq\lambda\bsrc(x_n)
\geq\lambda^{n+1}\bsrc(x_0)$.  Positivity follows from
$\lambda>0$ and $\bsrc(x_0)>0$.
\end{proof}

The theorem is conditional mathematics, not a derivation of $\Rsrc$,
$\Usrc$, $\bsrc$, or $\Osrc$.  Those provenance burdens are exactly what the
future Ontology Formalizer must resolve or return as a precise obstruction.

\subsection{Exact rational control}

On RT018's rank-two wall coordinate $t$, take $\Usrc=\{t>0\}$ and
$b(t)=t$.  A separately supplied transition $F_q(t)=qt$ with
$q\geq 1/3$ obeys
\[
  b(F_q(t))=qt\geq \tfrac13 b(t)>0.
\]
By contrast, the admitted map $t\mapsto -t$ crosses the wall from every
$t>0$.  Exact rational samples confirm both statements.  This fixture shows
what an independently sourced protection law would have to distinguish; it
does not supply that law and is not source evidence.

\section{The frozen classifier remains diagnostic only}

\begin{corollary}[Diagnostic nonselection]
The fixed $\Bridge$ readout may be evaluated along an independently selected
source trajectory, but its pointwise totality and presentation equivariance
cannot replace either $\Osrc$ or $\Rsrc$.
\end{corollary}

\begin{proof}
RT018 proves that pointwise totality selects no input, path, stable component,
or physical interpretation.  The occurrence--viability theorem shows that
selection and protection are separately required.  Evaluation of a
diagnostic after those structures are supplied does not make the diagnostic
their source.
\end{proof}

Accordingly, the selected packet may retain the full realization data and may
evaluate $\Bridge$ for diagnostics, but its selector may not factor only
through shared $\tau$, $Q_{\mathrm{ray}}$, common-character data, or the
unchanged $\Bridge$ components.  Presentation naturality remains mathematical
equivariance, not a physical gauge law.

\section{Four-route comparison}

\begin{longtable}{P{0.20\linewidth}P{0.30\linewidth}P{0.40\linewidth}}
\toprule
Route & Principal value & Disposition and reason \\
\midrule
\endfirsthead
\toprule
Route & Principal value & Disposition and reason \\
\midrule
\endhead
Source dynamical viability and admissibility theorem &
Targets occurrence selection and finite-variation protection through explicit
source domains, transitions, barriers, and failure branches. &
\status{Selected, not executed}.  It is constructive, materially independent
of all four freezes, target-free at the selector boundary, and admits a
decisive theorem-or-obstruction contract without protected authority. \\
Realization-sensitive margin bridge &
Retains quantitative moduli forgotten by named oriented-matroid components. &
\status{Not selected}.  A margin requires an independently sourced norm or
topology and still does not select which realization occurs.  Its added
structure and preload risk are higher than the selected route. \\
Source-side irrelevance theorem &
Could delimit a source-justified class of classifiers unable to determine
physical interpretation. &
\status{Not selected}.  RT018 already supplies the narrow unchanged-bridge
obstruction, no broader lawful bridge class is presently delimited, and a
constructive source-side packet remains available. \\
Protected human-gated ontology stop &
Correctly preserves authority if every honest continuation needs ontology
adoption or canonical modification. &
\status{Not selected}.  One target-free proposal-only formalization or precise
obstruction can proceed without adoption; a protected stop is premature. \\
\bottomrule
\end{longtable}

The numerical rubric in the task-local exact model is a declared routing aid,
not scientific evidence.  Its result agrees with the qualitative discriminator:
Route A uniquely addresses both occurrence and stability while avoiding the
unselected quantitative structure of Route B and the nonconstructive posture
of Routes C and D.

\section{Selected future packet contract}

The selected packet defines, or returns the exact obstruction to defining:
\begin{enumerate}
  \item a fully typed source-state or realization space $\Xsrc$;
  \item a source-provenanced transition relation or semiflow $\Rsrc$;
  \item a predeclared target-free admissible domain $\Usrc$;
  \item an exact barrier $\bsrc$ or an exact alternative viability
        certificate;
  \item the universal viability kernel $\Viab_{\Rsrc}(\Usrc)$;
  \item an independently source-provenanced occurrence or component selector
        $\Osrc$; and
  \item justified source covariance without promoting presentation arrows to
        physical gauge.
\end{enumerate}

The future completion must prove a theorem with explicit hypotheses and proof,
construct a complete proposal-only candidate law, or return the smallest
missing primitive or countermodel.  An obligation list without a decisive
result fails the packet.  Multiple viable components without an independent
selector fail closed.  Success-filtered admissibility, desired-cone filtering,
empirical calibration, target geometry, benchmark behavior, or imports from
the four frozen routes also fail closed.

\section{Source-extension and authority classification}

The current ontology does not derive the required occurrence selector,
admissible-transition relation, barrier propagation law, or unique viable
component.  The selected object is therefore classified as a
\status{proposal-only new ontology primitive candidate}, with
\status{blocked adoption and open continuation}.  This phrase does not assert
that the primitive is necessary in every possible extension or that its
future formalization will pass audit.

No target manifold, target atlas, target metric, tetrad, coframe, desired GR
cone, proper time, empirical detector response, physical normalization, or
benchmark answer is an admissible source premise.  Any such import triggers a
future smuggling-audit failure.  If a candidate law is formalized, the next
separately admitted role is a Smuggling Auditor; if that audit passes, Refuter
stress remains separately required.  There is no automatic adoption or B2
disposition.

\section{Distance-to-GR matrix}

\small
\begin{longtable}{P{0.23\linewidth}P{0.13\linewidth}P{0.55\linewidth}}
\toprule
Burden & Status & Exact reason \\
\midrule
\endfirsthead
\toprule
Burden & Status & Exact reason \\
\midrule
\endhead
Source ontology primitives & no delta & One proposal-only law packet is selected; no primitive is formalized or adopted.\\
Source equivalence $\mathrm{Eq}_{\mathrm{Src}}$ & no delta & No source equivalence or physical gauge law is established.\\
$\mathrm{RetainH}$ & no delta & No retention primitive or theorem is constructed.\\
$\mathrm{GenH}$ & no delta & No source generator or dynamics is constructed; $\Rsrc$ is a future burden.\\
$\mathrm{ObsLoc}_{\mathrm{lc}}$ & no delta & No observer-localization semantics is constructed.\\
$\mathrm{Resp}_{\mathrm{lc}}$ & no delta & Source viability and Bridge diagnostics are not empirical or physical response.\\
$M_{\mathrm{src}}$ & no delta & $\Xsrc$ is a future typed proposal domain, not an adopted smooth source manifold.\\
$g_{\mathrm{eff}}$ & no delta & No physical cone, conformal class, scale, source-to-metric map, or metric is derived.\\
Matter coupling & no delta & No universal sector propagation or coupling law is derived.\\
Einstein equations & no delta & No action, stress tensor, dynamics, correction law, or field equation is derived.\\
Finite-variation robustness & no delta & A conditional theorem target is selected, but no source protection law is formalized or proved for the current ontology.\\
Benchmark promotion & no delta & No target benchmark or GR recovery test is performed.\\
Gate Chair status & no delta & No protected review, adoption, activation, or verdict is requested.\\
Current route freeze or hard-fail status & no delta & All four local freezes persist; no global freeze or hard-fail verdict is added.\\
\bottomrule
\end{longtable}
\normalsize

The matrix contains fourteen literal no-delta rows.  Packet selection narrows
the next research action but does not discharge a derivation burden.

\section{Next route and non-conclusions}

After this transaction is checkpointed, exactly one separately admitted
Ontology Formalizer AgentJob may execute the selected packet.  The present
task does not execute it.  It does not reopen or supersede a freeze, correct a
predecessor, evaluate D7 or P2--T01 adequacy, activate, dispose, adopt, reject,
or certify B2, unlock P4--T03, derive physical causality or empirical response,
establish all-P7 sector coverage, construct conformal geometry or scale,
define $g_{\mathrm{eff}}$, derive matter coupling or Einstein equations,
request a protected Gate verdict, promote a benchmark or proof, authorize
publication, push, release, or external action, claim a global no-go, or claim
a completed derivation.

\section*{Control provenance}

This manuscript is the fused parent output of one outer AgentJob with two
internal physics perspectives.  Child agreement or conflict resolution,
source hashes, exact-model output, LaTeX compilation, validation, and a Git
checkpoint are control evidence only.  They do not establish the truth,
source provenance, physical interpretation, or adoption of the future law.

\end{document}
